\documentclass[14pt, a4paper]{extarticle}
\usepackage[utf8]{inputenc}						% Setting encoding
\usepackage[russian]{babel}						% Set russian landuage 
\usepackage{indentfirst}						
\usepackage{setspace}
\usepackage{geometry}
\usepackage{caption}							% Пакет для настройки плавующих элементов (Рисунки и т.д)
\usepackage{graphicx}
\usepackage{chngcntr}							% Для привязки нумерации русунка к сеуции
\usepackage{float}								% Для использования H для корректного отображения рисунков
\usepackage{url}
\usepackage{microtype}


\usepackage{tempora} 						 	% Times New Roman for text
\usepackage[bigdelims,vvarbb]{newtxmath}	 	% math italic letters from Nimbus Roman
\usepackage[cal=boondoxo]{mathalfa}  			% mathcal from STIX, unslanted a bit

\usepackage{longtable}
\usepackage{xltabular}
\usepackage{ragged2e}
\newcolumntype{L}{>{\RaggedRight\arraybackslash}X}

\babelfont[russian]{rm}[
Scale=MatchUppercase,
ItalicFont={Tempora-Italic.otf},
BoldFont={Tempora-Bold.otf},
BoldItalicFont={Tempora-BoldItalic.otf}
]{Tempora-Regular.otf}

% Рисунки
\counterwithin{figure}{section}					% Сама привязка рисунка

\captionsetup[figure]{
	labelsep=endash,								% Изменяем двоеточее на тире
	name={Рисунок}								% Изменение подписи рисунка
}				

\setlength{\intextsep}{1\baselineskip}  % Отступ вокруг figure

% Отступы
\geometry{left=3cm, right=1.5cm, top=2cm, bottom=2cm} % Отступы
\onehalfspacing
\setlength{\parindent}{1.25cm}
\setlength{\parskip}{0pt}


% Разделы
\makeatletter
\renewcommand\section{\@startsection{section}{1}{1.25cm}% 
	{-3.5ex \@plus - 1ex \@minus -.2ex}% 
	{2.3ex \@plus.2ex}%
	{
		\bfseries
		\fontsize{16}{1.5}\selectfont
	}}

\renewcommand{\thesection}{\arabic{section}}					% redefines the section numbering format 
\renewcommand{\@seccntformat}[1]{\csname the#1\endcsname\space}

\makeatother


\begin{document}
	\begin{titlepage}
		\centering
		Министерство науки и высшего образования Российской Федерации
		
		Федеральное государственное автономное образовательное учреждение высшего образования
		
		\MakeUppercase{Томский государственный университет систем управления и радиоэлектроники (ТУСУР)}
		
		\vspace{1\baselineskip}
		
		[Передовая инженерная школа (ПИШ)]
		
		\vspace{2\baselineskip}
		
		{\bfseries\MakeUppercase{
				
				[Название работы]
				
		}}
		
		Отчет по [тип работы] по дисциплине "[дисциплина]"
		
		\vspace{1\baselineskip}
		
		\begin{flushright}
			Студент(ы) гр. 362-2:
			
			{\rule{1.5cm}{0.6pt} И. О. Фамилия}
			
			
			"\rule{1.5cm}{0.6pt}" \rule{2cm}{0.6pt} {\the\year}г.
			
			\vspace*{1\baselineskip}
			
			Проверил:
			
			Должность (Доцент кафедры ПрЭ):
			
			{\rule{1.5cm}{0.6pt} И. О. Фамилия}
			
		\end{flushright}
		
		\vfill
		{Томск \the\year}
		
	\end{titlepage}
	
	\textbf{\MakeUppercase{Введение}}
	
	
	
	
	\clearpage
	\section{Основная часть}
	
	
	
	\clearpage
	
	\section{Примечания}
	
	
	
	\end{document}

